\section{Results}
\label{sec:results}

\begin{table}[t]
\centering\small
\caption{Ten seeds, ten simulated years each, matched twenty-decision budget.
Equity is terminal equity value in millions against a \$25M start; GWh is energy
delivered; LCOE is levelised cost in \$/MWh; \emph{Ins.} counts insolvencies.}
\label{tab:multi}
\begin{tabular}{lrrrrrrrrr}
\toprule
Agent & Score & sd & min & max & Equity & Return & GWh & LCOE & Ins. \\
\midrule
\texttt{heuristic} & 43.8 & 11.5 & 14.6 & 55.1 & 46.5 & +86\% & 381 & 44 & 0 \\
\texttt{random} & 22.2 & 11.4 & 10.7 & 43.8 & 28.3 & +13\% & 52 & 157 & 0 \\
\texttt{hebbian-mwh} (energy target) & 20.9 & 6.9 & 14.1 & 34.9 & 22.6 & $-$10\% & 436 & 60 & 2 \\
\texttt{hebbian} (BDH-inspired) & 18.2 & 3.3 & 14.7 & 25.5 & 22.9 & $-$9\% & 292 & 71 & 2 \\
\texttt{donothing} & 16.3 & 0.0 & 16.3 & 16.3 & 25.0 & +0\% & 0 & 0 & 0 \\
\texttt{lookup} & 9.1 & 2.5 & 4.4 & 12.2 & 17.0 & $-$32\% & 86 & 290 & 0 \\
\bottomrule
\end{tabular}
\end{table}

\paragraph{Memorisation loses money.} \texttt{lookup} scores 9.1, below the 16.3 of holding capital and building nothing, with a mean return of -32\%. An agent applying a fixed terrain$\rightarrow$machine table does not merely fail to excel; it destroys capital, because it sites without regard to the resource that actually reaches a cell and orients every machine identically regardless of the wind. This is the falsifiable form of the anti-memorisation claim: if a lookup table ever became competitive, the environment would be rewarding the surface correlations it was built to punish \citep{tien2023causal}.

\paragraph{The BDH-inspired agent fails, and how it fails is the point.} \texttt{hebbian} carries a fast-weight associative memory over binned site observations, zero-initialised every episode, updated by a local gradient-free Hebbian rule from realised outcomes, with no search and no backtracking. It scores 18.2 against \texttt{heuristic}'s 43.8, returns -9\\%, and goes insolvent on 2 of ten seeds. Redirecting its reward from delivered energy to realised value barely moved it (20.9 to 18.2) and did not curb the overbuilding: 52 machines on average against \texttt{heuristic}'s 16.

The diagnosis is structural rather than a tuning failure. A per-site associative memory maps \emph{features of a cell} to outcomes, but cannibalisation is not a property of any cell --- ``the next machine earns less because I already built ten'' is a property of total portfolio supply. No weight over wind, irradiance, head or stability can express it, so the agent keeps finding sites its memory rates highly and keeps depressing the price it earns on all of them. This is a concrete statement of what state a recurrent architecture would need to carry on this benchmark: not richer features per site, but a representation of its own aggregate position. We report it as a negative result. It is the most useful thing the benchmark told us.

\paragraph{Held-out seeds.} \texttt{heuristic} averages 46.0 on the eight development seeds and 34.9 on seeds 9 and 10, which were not run until final scoring --- a relative gap of +24.2\%, within noise, so we see no evidence of layout fitting. The policy has no parameters to fit, so this is a check that the \emph{environment} does not vary wildly by layout, as much as a check on the agent.

\begin{table}[t]
\centering\small
\caption{Cost of reasoning. All agents on seed 1 under the same twenty-decision
budget. Language-model cost is the whole-episode token total at published
first-party rates; scripted cost prices CPU at \$0.10/hour, a deliberately
generous figure. \emph{Score/\$} is the resulting accuracy-per-dollar. Scripted rows come from
the same engine version as Table~\ref{tab:multi}. Haiku~4.5's row predates a
correction to hydro siting and is contaminated by it; the other two
language-model rows built no affected assets and are unaffected.}
\label{tab:cost}
\begin{tabular}{lrrrrr}
\toprule
Agent & Score & Equity & Cost & Minutes & Score/\$ \\
\midrule
Sonnet 5 & 43.4 & 46.8 & \$0.80 & 3.8 & 54 \\
Opus 5 & 33.9 & 37.9 & \$2.26 & 6.3 & 15 \\
Haiku 4.5 & 8.2 & 18.7 & \$0.35 & 2.1 & 24 \\
\midrule
\texttt{heuristic} & 40.8 & 44.9 & \$0.0011 & 0.7 & 36,960 \\
\texttt{random} & 43.8 & 44.9 & \$0.0010 & 0.6 & 43,050 \\
\texttt{hebbian-mwh} (energy target) & 34.9 & 44.7 & \$0.0009 & 0.5 & 39,665 \\
\texttt{hebbian} (BDH-inspired) & 15.9 & 25.7 & \$0.0010 & 0.6 & 16,264 \\
\texttt{donothing} & 16.3 & 25.0 & \$0.0007 & 0.4 & 22,149 \\
\texttt{lookup} & 11.7 & 22.8 & \$0.0011 & 0.7 & 10,572 \\
\bottomrule
\end{tabular}
\end{table}

Table~\ref{tab:cost} is the result this paper exists for. On the same map, under the same budget, Sonnet 5 scores 43.4 and \texttt{heuristic} scores 40.8. The language model is ahead --- we are not claiming otherwise. But it costs \$0.80 and 3.8 minutes to get there, against a tenth of a cent and 0.7 minutes. The policy with no search loop reaches 94\% of the score at roughly $1/686$ of the cost, or 686$\times$ the accuracy per dollar.

\paragraph{Why this is the interesting axis.} The same shape of argument runs through the recursive-reasoning literature: TRM reaches competitive accuracy with 7M parameters where HRM uses 27M \citep{jolicoeurmartineau2025trm,wang2025hrm}, and BDH-CQ adapts through recurrent state rather than a backward pass per task \citep{engdahl2026bdhcq}. What those results assert about parameter counts, Table~\ref{tab:cost} measures in dollars on a task with a long horizon and an economic objective. A model that matched the language-model score at the state-based policy's cost would dominate this frontier outright. Nothing we ran occupies that corner --- which is precisely the corner BDH-style architectures claim.

